\chapter{Sphères, barycentres et mesures}
\section*{Exercice 8.1 -- Sphère et plan}
Le projeté du centre est $(0,0,3)$, à distance $3$. L'intersection est un cercle de rayon $\sqrt{25-9}=4$ dans ce plan. Son équation y est $x^2+y^2=16$.

Pour le second lieu, $MA^2=MB^2$ donne $x^2+y^2+z^2=(x-2)^2+y^2+z^2$, donc $x=1$. C’est le plan médiateur de $[AB]$, normal à $(1,0,0)$ et passant par le milieu $(1,0,0)$.

\section*{Exercice 8.2 -- Tangence}
Le rayon a pour direction $(0,0,1)$ ; le plan normal à ce vecteur passant par le point est $z=1$. L'équation de la sphère y devient $x^2+y^2=0$, donc le seul point commun est le point de tangence.

\section*{Exercice 8.3 -- Barycentre}
La somme des poids est $3\ne0$ et $G=(2A+B)/3=(1,0,0)$. Le point est plus proche du point portant le poids le plus grand ; cette observation contrôle le sens du résultat.

\section*{Exercice 8.4 -- Leibniz}
Appliquer l'identité de Leibniz avec poids $1,1$ : le terme constant vaut $IA^2+IB^2=2(AB/2)^2=AB^2/2$. La somme est minimale pour $M=I$.

\section*{Exercice 8.5 -- Volume}
La base dans $z=0$ est un triangle rectangle d'aire $2\times3/2=3$. La hauteur vaut $4$ car l'arête verticale est normale à ce plan. Le volume est $3\times4/3=4$ unités de volume.

\section*{Exercice 8.6 -- Sphère circonscrite}
Pour $C=(x,y,z)$, l'égalité des distances à $O$ et $(2,0,0)$ donne $x=1$ par soustraction des carrés. Les deux autres égalités donnent $y=z=1$. Le centre est $(1,1,1)$ et le rayon $\sqrt3$.
